% Source: https://artofproblemsolving.com/community/c3h1321747_geometry
% Stable identity: aops-topic-1321747-q01
Put the centre at $O=(0,0)$ and $A=(0,5)$. By symmetry $B=(m,y)$ and $C=(-m,y)$. Then $m^2+y^2=25$ and $AB^2=m^2+(5-y)^2=36$. Expanding the second equation, $m^2+y^2-10y+25=36$, and using the first, $50-10y=36$, so $y=\dfrac{7}{5}$, and $m^2=25-\dfrac{49}{25}=\dfrac{576}{25}$, $m=\dfrac{24}{5}$.

\begin{asy}
usepackage("amsmath");
size(4.5cm);
defaultpen(fontsize(10pt));
path ra(pair P,pair u,pair v,real s=0.2){return P+s*unit(u)--P+s*unit(u)+s*unit(v)--P+s*unit(v);}
pair O=(0,0); draw(circle(O,5),linewidth(0.9)); pair A=(0,5),B=(4.8,1.4),C=(-4.8,1.4);
draw(A--B--C--cycle,linewidth(0.9)); draw(O--B,gray+dashed); draw(A--O,gray+dashed); dot(O); label("$O$",O,S);
label("$A$",A,N);label("$B$",B,E);label("$C$",C,W); label("$6$",(A+B)/2,NE); label("$6$",(A+C)/2,NW);
pair F=(0,1.4); dot(F); label("$F$",F,SW); label("$5$",(O+B)/2,SE,gray);
label("\mbox{}",(min(currentpicture,user=true)+max(currentpicture,user=true))/2,black); // keeps the text after the figure black when LaTeX typesets the labels inline
\end{asy}

Hence $BC=2m=\boxed{9.6\text{ cm}}$. Since $y=\dfrac{7}{5}>0$, the centre lies below $BC$, outside the triangle, as stated (in fact this is forced by $AB=AC=6$). Alternatively, with $F$ the midpoint of $BC$: $AF=\dfrac{AB^2}{2R}=\dfrac{36}{10}=3.6$ and $BF=\sqrt{36-3.6^2}=4.8$.
